\phase{0 --- Architecture decisions and package boundaries}

\textbf{Objective.} Resolve the decisions that affect every protocol, implementation package, and application before freezing public API.

\textbf{Why first.} A mistake in operation identity, precision meaning, policy composition, or provider lifecycle propagates into every implementation. These are shared contracts, not implementation-specific details.

\deliverables
\begin{enumerate}
  \item A maintained Architecture Decision Register recording each foundational decision as baseline, proposed, or deferred, including rationale, invariants, consequences, and validation required before API freeze.
  \item Explicit decisions or tracked proposals for image/layer identity, resource ownership, snapshot/replacement concurrency, operation identity/versioning, and asynchronous completion.
  \item A formal boundary among \IPP{}, \GEP{}, \IPR{}, external provider packs, and external operation-implementation packs.
  \item A decision on typed public implementation protocols and private type erasure for the heterogeneous runtime registry.
  \item A specification schema distinguishing operation semantics, implementation support, provider static capabilities, and dynamic execution state.
  \item A precision vocabulary that distinguishes input/storage, computation, accumulation, and output precision.
  \item A qualification-policy model distinguishing hard requirements from soft preferences, with explicit scopes and deterministic composition rules.
  \item Conflict/error taxonomy: contradictory policy, unsupported semantic requirement, missing provider, temporary device unavailability, resource reservation failure, and implementation execution failure.
  \item A platform/toolchain support baseline and an environment telemetry capability matrix. Record which state metrics are reliable, optional, estimated, or unavailable.
  \item Repository/package boundaries, dependency-cycle checks, command-line build/test entry points, and CI for the initially supported targets.
\end{enumerate}

\exitcriteria
\begin{itemize}
  \item No unresolved decision blocks the first protocol implementation.
  \item The dependency graph is acyclic: \IPP{} may depend on \GEP{}; the reverse dependency is forbidden.
  \item No core package depends on a concrete CPU/GPU/algorithm implementation package.
  \item The policy model can represent the required FP32 example without confusing storage and computation precision.
  \item The design distinguishes contradictory policy from a valid requirement that cannot currently be satisfied by available hardware.
  \item No foundational decision is left implicit: unresolved items have an owner/status and do not block the work that can safely proceed.
\end{itemize}

\textbf{Do not do yet:} write provider kernels, add a large operation catalog, or create a general-purpose dynamic plug-in loader.

\phase{1 --- Image operation and implementation contracts}

\textbf{Objective.} Specify the domain contracts that allow one operation to have many independent implementation classes.

\deliverables
\begin{enumerate}
  \item \textbf{Operation identity:} stable identifier, semantic version, immutable parameters, parameter validation, and typed input/output contracts.
  \item \textbf{Operation specification:} semantic requirements, permitted numerical variation, supported input/output categories, parameter domain, mutation/effect semantics, and precision dimensions.
  \item \textbf{One generic implementation protocol:} \texttt{ImageOperationImplementation} with an associated \texttt{Operation} type, declarative specification, and typed execution entry point.
  \item \textbf{Implementation specification:} operation identity/version, implementation identity/version, provider/runtime dependency, required generic capabilities, supported formats/layouts/precision modes, parameter limits, semantic options, resource requirements, and static cost hints.
  \item \textbf{No per-implementation qualifier:} implementations declare support; they do not each implement independent \texttt{qualify()} code. The core resolver interprets their specifications.
  \item \textbf{Qualification result vocabulary:} eligible, rejected, and temporarily unavailable outcomes with machine-readable reasons, requirements, and preparation/transfer details.
  \item \textbf{Fluent policy contract:} immutable pipeline/operation-scoped configuration with explicit \texttt{require} and \texttt{prefer} semantics.
  \item \textbf{Example contracts:} Gaussian blur, resize, and blending with enough detail to compare CPU, CUDA, Metal, and algorithm-library implementations without changing operation identity.
\end{enumerate}

\exitcriteria
\begin{itemize}
  \item Two different implementation classes can declare support for the same \texttt{GaussianBlur} operation without defining separate operation types.
  \item A third implementation can be added without modifying the operation type or generic implementation protocol.
  \item Static implementation capability is distinguishable from runtime provider availability and transient load.
  \item FP32 computation can be expressed separately from FP32 input/output storage and accumulation precision.
  \item No public contract requires provider-native types or per-backend qualification functions.
\end{itemize}

\textbf{Gate.} Do not begin concrete algorithm packages until an operation's semantics and the implementation-specification schema can express meaningful support differences.

\phase{2 --- GeneralExecutionProtocol}

\textbf{Objective.} Define generic execution resources, provider lifecycle, static capabilities, and dynamic state observations without introducing image semantics.

\textbf{Dependency constraint.} \GEP{} MUST NOT import \IPP{}, define Image/Layer/Pixel semantics, or know operation names such as blur, blend, resize, or segmentation.

\deliverables
\begin{enumerate}
  \item Provider identity/version and platform-aware provider discovery.
  \item Device identity, properties, supported features, execution contexts, and static capability descriptors.
  \item Generic resource identifiers scoped by execution target, descriptors for buffer/texture/opaque resource kinds, capacity/alignment/memory-domain metadata, and validated allocation requests.
  \item Optional allocation/import context capabilities. Imported handles retain their source resources; no image/pixel semantics enter \GEP{}.
  \item Resource access declarations for read, write, and read/write use, with reference-based lifetime semantics.
  \item Typed work submission with declared resource accesses and completion dependencies; providers retain work and resources through completion.
  \item Typed submission/completion handles, cancellation outcomes, dependency-failure behavior, and optional context/resource synchronization capabilities.
  \item Submission contracts that preserve native differences: Metal command buffers, CUDA streams, OpenCL queues, and D3D command lists remain provider-owned rather than being forced into one public native-command object.
  \item A state/availability observation model for device health, availability, memory pressure, outstanding submissions, queue depth, and estimated delay where available.
  \item Observation timestamps and freshness/quality metadata. Unsupported telemetry is represented as unknown, never as zero load.
  \item Structured provider errors, device-loss behavior, allocation/retention, and release rules.
  \item Fake/test provider and conformance doubles whose capabilities and state can be controlled deterministically.
\end{enumerate}

\exitcriteria
\begin{itemize}
  \item \GEP{} builds and tests without referencing image-processing symbols.
  \item Resource request tests verify kind, capacity, alignment, memory domain, and context target.
  \item Imported resource handles retain the source allocation for the full import lifetime.
  \item A submitting provider retains the work item and declared resources until completion, even if caller references are dropped.
  \item Dependencies complete successfully before dependent work starts; failed dependencies do not silently execute dependent work.
  \item Typed submission output is valid after successful completion. Cancellation and provider errors preserve resource lifetime until work is no longer in use.
  \item Optional synchronization is scoped and does not require every provider to expose the same native primitive.
  \item Native provider differences remain behind optional/provider-specific capabilities; no universal command-buffer fiction is introduced.
  \item Stale or unknown state is never misreported as current or idle.
\end{itemize}

\textbf{Gate.} Do not make the common protocol larger merely to expose one vendor-specific metric or command primitive.